% !TeX root = ../main.tex
% ============================================================
% LÁMINA 13 — ETAPA 6: FHIR (fondo claro)
% Datos: latex/capitulos/06_Resultados_y_Discusion.tex (Bundle FHIR, Objetivo 4)
%        latex/capitulos/07_Conclusiones.tex (cumplimiento Objetivo 4)
%        latex/tablas/capitulo_evaluacion/tabla_e2e_detalle_{batch,streaming}.tex
% ============================================================
\begin{frame}[plain]
  \begin{tikzpicture}[remember picture, overlay]

    \fill[paperBg] (current page.south west) rectangle (current page.north east);

    \node[anchor=north west, align=left]
      at ([xshift=1.4cm, yshift=-0.55cm]current page.north west)
      {\kicker[accentBlue]{12 \,$\cdot$\, Caso clínico \,$\cdot$\, Etapa 6}};

    \node[anchor=north west, align=left, text width=13.2cm]
      at ([xshift=1.4cm, yshift=-1.05cm]current page.north west)
      {\large\bfseries\sffamily\color{inkText}El \textit{Bundle} \textit{FHIR} del caso: válido en forma, no en contenido};

    % ------------------------------------------------------
    % Breadcrumb de 6 etapas (Validación resaltada)
    % ------------------------------------------------------
    \def\bcW{2.05cm}
    \def\bcH{0.5cm}
    \def\bcStep{2.21cm}
    \def\bcY{-1.70cm}

    \node[anchor=north west, fill=cardWhite, draw=lineGrey, line width=0.6pt, rounded corners=8pt,
          minimum width=\bcW, minimum height=\bcH, align=center, inner sep=0pt]
      at ([xshift=1.4cm+0*\bcStep, yshift=\bcY]current page.north west)
      {\fontsize{6.5}{6.5}\selectfont\bfseries\sffamily\color{mutedText}1\,\textperiodcentered\,Captura};

    \node[anchor=north west, fill=cardWhite, draw=lineGrey, line width=0.6pt, rounded corners=8pt,
          minimum width=\bcW, minimum height=\bcH, align=center, inner sep=0pt]
      at ([xshift=1.4cm+1*\bcStep, yshift=\bcY]current page.north west)
      {\fontsize{6.5}{6.5}\selectfont\bfseries\sffamily\color{mutedText}2\,\textperiodcentered\,ASR};

    \node[anchor=north west, fill=cardWhite, draw=lineGrey, line width=0.6pt, rounded corners=8pt,
          minimum width=\bcW, minimum height=\bcH, align=center, inner sep=0pt]
      at ([xshift=1.4cm+2*\bcStep, yshift=\bcY]current page.north west)
      {\fontsize{6.5}{6.5}\selectfont\bfseries\sffamily\color{mutedText}3\,\textperiodcentered\,NER};

    \node[anchor=north west, fill=cardWhite, draw=lineGrey, line width=0.6pt, rounded corners=8pt,
          minimum width=\bcW, minimum height=\bcH, align=center, inner sep=0pt]
      at ([xshift=1.4cm+3*\bcStep, yshift=\bcY]current page.north west)
      {\fontsize{6.5}{6.5}\selectfont\bfseries\sffamily\color{mutedText}4\,\textperiodcentered\,Normaliz.};

    \node[anchor=north west, fill=cardWhite, draw=lineGrey, line width=0.6pt, rounded corners=8pt,
          minimum width=\bcW, minimum height=\bcH, align=center, inner sep=0pt]
      at ([xshift=1.4cm+4*\bcStep, yshift=\bcY]current page.north west)
      {\fontsize{6.5}{6.5}\selectfont\bfseries\sffamily\color{mutedText}5\,\textperiodcentered\,Generación};

    \node[anchor=north west, fill=inkDark, rounded corners=8pt,
          minimum width=\bcW, minimum height=\bcH, align=center, inner sep=0pt]
      at ([xshift=1.4cm+5*\bcStep, yshift=\bcY]current page.north west)
      {\fontsize{6.5}{6.5}\selectfont\bfseries\sffamily\color{accentTeal}6\,\textperiodcentered\,Validación};

    % ------------------------------------------------------
    % Izquierda: validación estructural superada
    % Derecha: advertencia de corrección semántica no garantizada
    % ------------------------------------------------------
    \node[anchor=north west, fill=cardWhite, draw=lineGrey, line width=0.6pt,
          rounded corners=6pt, minimum width=5.9cm, minimum height=4.35cm,
          align=left, inner sep=0pt]
      (fhirbox) at ([xshift=1.4cm, yshift=-2.45cm]current page.north west) {};
    \fill[accentTeal, rounded corners=2pt]
      (fhirbox.north west) rectangle ([yshift=-0.12cm]fhirbox.north east);
    \node[anchor=north west]
      at ([xshift=0.4cm, yshift=-0.35cm]fhirbox.north west)
      {\footnotesize\sffamily\color{mutedText}\MakeUppercase{Validación estructural}};
    \node[anchor=north west]
      at ([xshift=0.4cm, yshift=-1cm]fhirbox.north west)
      {\Large\bfseries\sffamily\color{accentTeal}FHIR válido};
    \node[anchor=north west]
      at ([xshift=0.4cm, yshift=-1.75cm]fhirbox.north west)
      {\parbox[t]{5.1cm}{\raggedright\fontsize{7.5}{9.2}\selectfont\sffamily\color{inkText}%
       El \textit{JSON FHIR} de este caso supera la validación estructural: \texttt{resourceType},
       referencias \texttt{subject} y \texttt{meta.profile} correctos.\\[4pt]
       }};

    \node[anchor=north west, fill=cardWhite, draw=lineGrey, line width=0.6pt,
          rounded corners=6pt, minimum width=6.4cm, minimum height=4.35cm,
          align=left, inner sep=0pt]
      (warnbox) at ([xshift=7.6cm, yshift=-2.45cm]current page.north west) {};
    \fill[accentRed, rounded corners=2pt]
      (warnbox.north west) rectangle ([yshift=-0.12cm]warnbox.north east);
    \node[anchor=north west]
      at ([xshift=0.4cm, yshift=-0.35cm]warnbox.north west)
      {\footnotesize\sffamily\color{mutedText}\MakeUppercase{Advertencia semántica}};
    \node[anchor=north west]
      at ([xshift=0.4cm, yshift=-1cm]warnbox.north west)
      {\parbox[t]{5.6cm}{\raggedright\fontsize{7}{8.4}\selectfont\sffamily\color{inkText}%
       La validación \textbf{no filtra} el contenido: el \textit{Bundle} traslada, sin corregirlos,
       los errores ya vistos en este mismo caso.\\[3pt]
      %  \textbf{\color{accentRed}NER:} \textit{"analgesia"} persiste como \texttt{Condition} (falso positivo).\\[2pt]
       \textbf{\color{accentRed}Terminología:} \textit{"doloroso"} $\to$ \textit{Trigeminal neuralgia}
       codificado como hallazgo (\textit{CIE-10} err.).\\[2pt]
       \textbf{\color{accentRed}Informe:} dosis $\times$1000 (\textit{"50\,mg/hora"}) pasa intacta al
       \texttt{MedicationStatement}.}};

    % ------------------------------------------------------
    % 3 tarjetas de estadísticas
    % ------------------------------------------------------
    % \def\statCardW{4.2cm}
    % \def\statCardH{1.6cm}
    % \def\statCardGap{0.3cm}
    % \def\statCardY{-6.95cm}

    % \node[anchor=north west, fill=inkDark, rounded corners=6pt,
    %       minimum width=\statCardW, minimum height=\statCardH, inner sep=0pt]
    %   (st1) at ([xshift=1.4cm, yshift=\statCardY]current page.north west) {};
    % \node[anchor=north west, align=left]
    %   at ([xshift=0.4cm, yshift=-0.10cm]st1.north west)
    %   {\Large\bfseries\sffamily\color{accentTeal}4\,/\,5};
    % \node[anchor=north west]
    %   at ([xshift=0.4cm, yshift=-0.70cm]st1.north west)
    %   {\parbox[t]{3.4cm}{\raggedright\fontsize{8}{9.5}\selectfont\sffamily\color{white}Puntuación \textit{FHIR} de este caso (\textit{streaming} Test\,1)}};

    % \node[anchor=north west, fill=inkDark, rounded corners=6pt,
    %       minimum width=\statCardW, minimum height=\statCardH, inner sep=0pt]
    %   (st2) at ([xshift=1.4cm+\statCardW+\statCardGap, yshift=\statCardY]current page.north west) {};
    % \node[anchor=north west, align=left]
    %   at ([xshift=0.4cm, yshift=-0.10cm]st2.north west)
    %   {\Large\bfseries\sffamily\color{accentTeal}3};
    % \node[anchor=north west]
    %   at ([xshift=0.4cm, yshift=-0.70cm]st2.north west)
    %   {\parbox[t]{3.4cm}{\raggedright\fontsize{8}{9.5}\selectfont\sffamily\color{white}Persisten como hallazgos clínicos válidos}};

    % \node[anchor=north west, fill=inkDark, rounded corners=6pt,
    %       minimum width=\statCardW, minimum height=\statCardH, inner sep=0pt]
    %   (st3) at ([xshift=1.4cm+2*\statCardW+2*\statCardGap, yshift=\statCardY]current page.north west) {};
    % \node[anchor=north west, align=left]
    %   at ([xshift=0.4cm, yshift=-0.10cm]st3.north west)
    %   {\Large\bfseries\sffamily\color{accentTeal}100\,\%};
    % \node[anchor=north west]
    %   at ([xshift=0.4cm, yshift=-0.70cm]st3.north west)
    %   {\parbox[t]{3.4cm}{\raggedright\fontsize{8}{9.5}\selectfont\sffamily\color{white}Objetivo 4 solo en forma}};

  \end{tikzpicture}
\end{frame}
